\chapter*{Sammendrag}

Den økende bruken av digitale bilder og video, sammen med strengere og stadig mer omfattende personvernregler, har ført til et større behov for pålitelige metoder for å beskytte sensitiv informasjon. Manuell sladding er tidkrevende og vanskelig å skalere, mens helautomatiske løsninger ofte har varierende nøyaktighet i komplekse visuelle omgivelser. Denne oppgaven presenterer \gsr{SafeVision}, et system for å oppdage og anonymisere sensitivt innhold i visuelle medier.

Systemet kombinerer objektdeteksjon og tekstanalyse for å identifisere sensitivt innhold i visuelle medier. Det kan oppdage ansikter, tekst og generelle objekter som personer og kjøretøy. I tillegg er det trent opp og integrert egne modeller i systemet for å klassifisere og merke domenespesifikke entiteter.

Systemet er utformet som en skalerbar og brukerorientert arkitektur som legger til rette for pålitelig sladding. Maskinlæringskomponentene er integrert i en Python basert backend, utviklet med FastAPI for asynkron behandling, sammen med en webbasert frontend utviklet med React og TypeScript. Brukere kan gjennomgå og justere deteksjoner før sladding, noe som muliggjør brukerstyrt forbedring og bidrar til økt nøyaktighet og bedre brukervennlighet. Den modulære arkitekturen, med tydelig separasjon mellom komponenter, gjør systemet enklere å videreutvikle og vedlikeholde, samtidig som lagring på klientsiden reduserer datatransport og styrker personvernet.